\section{问题重述}

\subsection{问题背景}
% 用自己的话概述题面背景（不要把题目原文整段抄下来）。篇幅控制在半页以内。

\subsection{问题提出}
% 逐条列出题面明确编号的顶层问题。只列编号问题，不要把背景描述、数据说明、提交要求
% 误当成子问题。子问题数按题面实际，不硬凑三个。

\textbf{问题 1：}

\textbf{问题 2：}

\textbf{问题 3：}

% 图 1：分析流程图（`fig_roadmap`）排在这里 —— 问题重述之后、问题分析之前。
%   参考件 `1_restatement.tex` 第 40 行就是这一句，
%   一字不差：`\paperfigure[0.85\textwidth]{fig_roadmap}{分析流程图}{fig_roadmap}`。
%   全文唯一的大图、唯一带配色的图；其余流程图/示意图一律半栏以内、黑白。
%   不得挪进 2_analysis.tex；**别缩，0.85 是下限** —— 954×1297 的竖版大图缩到
%   七成栏宽，节点字直接糊掉。

% 下面这句**必须取消注释**。模板里留成注释，是因为模板目录下没有 figures/*.pdf，
%    放开会让模板本身编译不过（`! Unable to load picture`）。
%    `python lib/web/check_skeleton.py` 会核它有没有放开 —— 放开是硬要求，别忘。
% \paperfigure[0.85\textwidth]{fig_roadmap}{分析流程图}{fig_roadmap}
